\documentclass[slide05.tex]{subfiles}
\begin{document}
\begin{frame}[plain]
    \begin{center}
        \vspace{48pt}
        {\huge\bf 演習}
    \end{center}
\end{frame}

\begin{frame}[fragile]
    \frametitle{Faboや赤外線でできることを考えよう}
    \begin{center}
        \begin{itemize}
            \item 例えば\dots
            \begin{itemize}
                \item ボリュームで点くLEDを切り替える
                \item 振動子がリズムよく震える
                \item OLEDでタイマーを作ってみる
                \item 距離センサを使って人が近くに来たら扇風機を回す
            \end{itemize}
            \item 例を参考に何が作れるか自分で考えてみよう
        \end{itemize}
    \end{center}
\end{frame}

\begin{frame}[fragile]
    \frametitle{作りたいものを考えよう}
    \begin{center}
        \begin{itemize}
            \item Faboや赤外線を使って作りたいものを考えよう(5〜10分)
            \item チームのみんなで作りたいものを発表して、一緒に考えてみよう(5〜10分)
            \begin{itemize}
                \item むずかしさはどれくらいだろう
                \item どうやって作ればいいだろう
            \end{itemize}
        \end{itemize}
    \end{center}
\end{frame}

\begin{frame}[fragile]
    \frametitle{手順}
    \begin{center}
        \begin{enumerate}
            \item 演習用のフォルダを作る
            \item プログラミングしてみよう
            \begin{itemize}
                \item hspファイルを書き換えてみる
                \item 新しくhspファイルを作ってみる
            \end{itemize}
        \end{enumerate}
    \end{center}
\end{frame}

\begin{frame}[fragile]
    \frametitle{できたものをチームのみんなに教えよう}
    \begin{center}
        \begin{itemize}
            \item どんなものを作ったか
            \item 工夫したところはどこか
            \item どうしてその作品を作ったのか
        \end{itemize}
    \end{center}
    {実際に作品を動かしてみんなに見せてみよう！}
\end{frame}
\end{document}
